%凍結延長されたらここを編集する
\newcommand{\TouketsuDate}{令和 11 年 3 月 31 日}


\section{特別法人税}

\subsection{課税の根拠}

\begin{shortanswer}[年金1 2020 問題2(4)\textcircled{1}]
  特別法人税が徴収される理由（創設当時における考え方）を簡記しなさい。
\end{shortanswer}

\begin{answerbox}
  事業主が拠出した掛金は、課税すべき従業員が特定されないことから、従業員へのみなし給与と
して拠出時に直ちに課税する代わりに、実際に年金の支給を受ける際に課税することとし、それ
までの課税繰延の利益を総体的に法人税として徴収するという考え方が取られた。
\end{answerbox}

\begin{shortanswer}[年金1 2020 問題2(4)\textcircled{2}]
特別法人税は1999 年（平成 11 年）4 月より課税が凍結されているが、その背景を、現在延
長されている具体的な凍結期間と合わせて簡記しなさい。
\end{shortanswer}

\begin{answerbox}
  バブル崩壊以降の低金利の継続や企業年金の財政状況を踏まえて、1999 年（平成11 年）4 月より
課税が凍結され、その後、課税凍結の延長が繰り返され、現在は \TouketsuDate 
まで課税が凍結されている。
\end{answerbox}

\begin{shortanswer}[年金1 2023 問題2(5)ア]
 特別法人税を課税することの考え方について、簡潔に入力しなさい。 （250 字以内）
\end{shortanswer}

\begin{answerbox}
特別法人税は、掛金拠出時に給与所得として課税せず、
給付時まで課税を繰り延べられることを踏まえ、
その期間の遅延利息相当分を課税するものである。
\end{answerbox}

\begin{shortanswer}[年金1 2023 問題2(5)ア]
特別法人税は、超低金利の状況等を踏まえて1999 年度（平成11 年度）から凍結されてはい
るが、確定給付企業年金、確定拠出年金の年金資産のうち、それぞれの課税対象部分につい
て、簡潔に入力しなさい。なお、具体的な税率については触れなくてよい。 （250 字以内）
\end{shortanswer}

\begin{answerbox}
  ＜確定給付企業年金＞

  年金資産のうち、加入者拠出分を除いた部分

  ＜確定拠出年金＞

  年金資産の全額
\end{answerbox}

\newpage

\subsection{税率}

\begin{fillblank}[数理人試験 2024 問題6 設問4]
  次は、退職年金等積立金に対する現時点の法人税の税率(1.0\%)の算定根拠である。
  A〜Dの空欄に入る数値をそれぞれ記載せよ。
  なお、A、Bは整数で、Cは小数点以下を切り捨てた整数で、Dは小数点以下第3位までで解答せよ。

  \begin{flalign*}
    & \{ （　　\mathrm{A}　　）\% +  （　　\mathrm{B}　　）\% \} \times （　　\mathrm{C}　　）\% \times \dfrac{1}{（　　\mathrm{D}　　）} \fallingdotseq 1.0\% \\
    & \Biggl\{ \begin{array}{l}
        \text{給与所得者の}\\
        \text{所得税の平均}\\
        \text{上積税率}
    \end{array} + \text{住民税率}\Biggr\} \times 
    \begin{array}{l}
        \text{日歩2銭の}\\
        \text{利子税率}
    \end{array} \times
    \begin{array}{l}
        \text{法人住民税と}\\
        \text{法人税の割合}
    \end{array} 
  \end{flalign*}
\end{fillblank}

\begin{answerbox}
  \begin{enumerate}[A:]
    \item 12
    \item 5
    \item 7
    \item 1.173
  \end{enumerate}
\end{answerbox}

\begin{truefalse}[年金1 H29 問題1(6)\textcircled{2}（改題）]
      \textbf{下線部の内容の正誤を判定し、誤っている場合は正しい内容に改めなさい。}
\\  

  特別法人税は、退職年金等積立金の額（確定給付企業年金は加入者負担分を除いた部分）に
対して国税として年率 \underline{1.173％}を乗じて計算した額、さらにこれに地方税（住民税）が課される
こととなっているが、現在は \TouketsuDate まで課税が凍結されている。
\end{truefalse}

\begin{answerbox}
  \textbf{誤}

  1.000\%
\end{answerbox}

\begin{explanationbox}
  特別法人税率の内訳：$法人税(1.000\%)+地方法人税(0.107\%)+法人住民税(0.070\%)$
\end{explanationbox}

\begin{fillblank}[数理人試験 H19 問題1(9)（改題）]
  確定給付企業年金・確定拠出年金の積立金には特別法人税が課税される。
  税率は、積立金（確定給付企業年金については（　　A　　）分を控除後）に対して（　　B　　）％（地方税分と合計で（　　C　　）％）である。
  なお、現在は（　　D　　）の規定により特別法人税は課税停止されている。

  厚生年金基金の積立金は（　　E　　）である。
  ただし、一定の給付水準を越える場合には、超過部分に相当する積立金に対し確定給付企業年金・確定拠出年金と
  同様に特別法人税が課税される。なお、課税判定の指標となる一定の給付水準とは、
  代行部分の（　　F　　）倍である。
\end{fillblank}

\begin{answerbox}
  \begin{enumerate}[A:]
    \item 加入者拠出
    \item 1.000
    \item 1.173
    \item 租税特別措置法
    \item 非課税
    \item 3.23
  \end{enumerate}
\end{answerbox}

\newpage

\subsection{その他}

\begin{truefalse}[年金1 2018 問題2(3)\textcircled{2}]
      \textbf{以下の文章の内容の正誤を判定し、誤っている場合は正しい内容に改めなさい。}
\\  

確定給付企業年金において加入者が負担した掛金は、拠出時に生命保険料控除の対象とな
り、運用時は特別法人税の課税対象とならない。
\end{truefalse}

\begin{answerbox}
  \textbf{正}
\end{answerbox}

\begin{truefalse}[年金1 H29 問題1(6)\textcircled{1}（改題）]
      \textbf{下線部の内容の正誤を判定し、誤っている場合は正しい内容に改めなさい。}
\\  

バブル崩壊以降の低金利の継続や企業年金の財政状況を踏まえて、 \underline{平成 14 年 4 月}から特別法
人税の課税は凍結され、その後、課税凍結の延長が繰り返され、現在は \TouketsuDate ま
で凍結されている。
\end{truefalse}

\begin{answerbox}
  \textbf{誤}

  平成11年4月
\end{answerbox}

\newpage